\section{Introduction}

In positive characteristic, equality of the dominant degrees in a
Frobenius coincidence problem can conceal a substantial cancellation.
For the H\'enon automorphism
\begin{equation}\label{eq:H-intro}
 H(x,y)=(y,y^q+g(y)-ax)
\end{equation}
over $\F_q$, both $H^n$ and the $n$th iterate of $q$-power Frobenius
have degree $q^n$. In their second coordinates the dominant monomials
cancel. The remaining degree, rather than the original degree, controls
the affine intersection. A calculation at periods prime to the
characteristic does not settle this problem: additional cancellations
appear when the period acquires successive factors of $p$.

We compute the remaining degree for all periods and all coefficients in
the following family. Throughout the paper,
\begin{equation}\label{eq:assumptions}
 q=p^e,\quad K=\overline{\F}_q,\quad
 a\in\F_q^*,\quad g\in\F_q[y],\quad
 2\le m:=\deg g<q,\quad p\nmid m.
\end{equation}
The leading coefficient of $g$ is denoted by $b\in\F_q^*$; no condition
is imposed on its other coefficients. Set
\[
 \Phi(x,y)=(x^q,y^q),\qquad S=H^{-1}\circ\Phi,
 \qquad E_n=\{H^n(P)=\Phi^n(P)\}\subset\A^2_K.
\]
Here $E_n$ denotes the closed equalizer scheme. Its geometric points are
the fixed points of $S^n$, as will also hold scheme-theoretically.

\begin{theorem}[All-period resonant count]\label{thm:count}
Assume \eqref{eq:assumptions}. For every $n\ge1$, the scheme $E_n$ is
finite and reduced. Write $r=p^{v_p(n)}$ and $s=n/r$. Then
\begin{equation}\label{eq:count}
 N_n:=\#E_n(K)
 =\frac{(m-1)q^{2n}+(q-m)q^{2n-r}}{q-1}.
\end{equation}
More precisely, the two literal affine equations of $E_n$ have respective
unique highest homogeneous terms
\begin{equation}\label{eq:top-intro}
 -x^{q^n},\qquad s b^r\overline m^{\,r-1}y^{d_n},
 \qquad
 d_n=q^{n-r}\frac{(m-1)q^r+(q-m)}{q-1},
\end{equation}
where $\overline m\in\F_p^*$ is the residue of $m$.
Their actual degrees are $q^n$ and $d_n$.
\end{theorem}

The expression for $d_n$ is an integer by the recurrence proved below;
its displayed denominator is an integer denominator, not division in
the coefficient field. In contrast, the coefficient $s$ in
\eqref{eq:top-intro} is read in $\F_p^*$ and survives because $p\nmid s$.
This distinction separates the arithmetic of degrees from the
characteristic-$p$ cancellation of coefficients.

The exact count also determines the analytic behavior of the associated
Artin--Mazur dynamical zeta germ
\begin{equation}\label{eq:zeta-definition}
 Z_S(t)=\exp\left(\sum_{n\ge1}N_n\frac{t^n}{n}\right).
\end{equation}

\begin{corollary}[Product and natural boundary]\label{cor:zeta}
Put $B=(q-m)/(q-1)$ and $u=q^2t$. With each logarithm normalized to vanish
at zero,
\begin{equation}\label{eq:product-intro}
 Z_S(u/q^2)=(1-u)^{-m/q}
   \prod_{k\ge1}(1-u^{p^k})^{e_k},\qquad
 e_k=\frac{B}{p^k}\bigl(q^{-p^{k-1}}-q^{-p^k}\bigr)>0.
\end{equation}
The product converges locally uniformly on $|u|<1$. For every fixed
integer $M\ge1$, the circle $|t|=q^{-2}$ is a natural boundary for
meromorphic continuation of $Z_S(t)^M$. In particular, $Z_S$ is neither
rational nor algebraic over $\C(t)$.
\end{corollary}

The algebraic point of the paper is the all-period cancellation formula
\eqref{eq:top-intro}, not a new general principle relating distorted
counts to natural boundaries. That principle has important antecedents.
Byszewski, Cornelissen, and Houben study endomorphisms of algebraic
groups with finite fixed sets at every iterate, develop finite-adelically
distorted counting sequences, and prove corresponding zeta dichotomies
\cite[Theorems A and C, and \S5.2]{bch2024groups}.
Their vector-group theorem applies to additive group endomorphisms,
whereas \eqref{eq:H-intro} yields a genuinely nonlinear affine-plane
problem. Proposition~\ref{prop:no-vector} gives a fixed-cardinality
obstruction to direct vector-group conjugacy in any dimension.
It does not exclude every possible group quotient presentation.

One-dimensional polynomial dynamics also has well-established
positive-characteristic zeta phenomena. In the version of Bridy's paper
specified in the bibliography, Theorem~1 gives rationality when the
derivative of a polynomial on $\A^1$ is zero, and Theorem~2 gives
transcendental examples \cite{bridy2012transcendence}.
These statements concern one variable, so neither determines the counts
of the two-dimensional map $S$. We make no claim to the first
positive-characteristic transcendental dynamical zeta function.

Quotients require a separate distinction. Theorems~A and B of the
specified preprint of \cite{bch2020affine} concern dynamically affine maps
on $\mathbb P^1$ and specified Kummer varieties; that work also supplies
a higher-dimensional framework under additional hypotheses.
We neither construct such a quotient for $S$ nor verify those hypotheses
for it. The present proof instead works directly in the polynomial ring.
The source comparison identifies the nonlinear leading-degree calculation
as the specific contribution considered here; it is a bounded comparison,
not a certificate of global priority.

Section~\ref{sec:system} identifies the single dynamical system and its
pullbacks. Section~\ref{sec:leading} resolves the leading terms at all
periods. Section~\ref{sec:count} proves finiteness, length, and reducedness.
Section~\ref{sec:zeta} derives the product and proves its boundary
statement directly. Section~\ref{sec:boundaries} records precise scope
and bounded exact checks.
